\documentclass[a4paper,12pt]{article}
\usepackage[T2A]{fontenc}
\usepackage[utf8]{inputenc}
\usepackage[russian]{babel}
\usepackage{amsmath,amssymb,icomma}
\usepackage[a4paper,margin=20mm,headheight=25pt]{geometry}
\usepackage{microtype,booktabs,tabularx,array,xcolor,tikz,fancyhdr,hyperref}
\definecolor{accent}{HTML}{008080}
\definecolor{darkaccent}{HTML}{004D4D}
\definecolor{err}{HTML}{9A302B}
\hypersetup{colorlinks=true,linkcolor=darkaccent,urlcolor=accent}
\setlength{\parindent}{0pt}
\setlength{\parskip}{5pt}
\setlength{\headsep}{7mm}
\pagestyle{fancy}
\fancyhf{}
\fancyhead[L]{\small ОГЭ. Вариант 21}
\fancyhead[R]{\small Ярослав Верещагин}
\fancyfoot[C]{\small\thepage}
\renewcommand{\headrule}{\hbox to\headwidth{\color{accent}\leaders\hrule height .5pt\hfill}}
\newcommand{\signature}{Лёвин Артём Александрович\\эксклюзивно для Ярослава Верещагина}
\newcommand{\task}[2]{\clearpage\phantomsection\label{task:#1}{\Large\color{darkaccent} Задание №#1. #2}\par\vspace{2mm}{\color{accent}\hrule}\vspace{2mm}}
\newcommand{\sub}[1]{\par\smallskip\textcolor{darkaccent}{\textbf{#1}}\par}
\newcommand{\error}[1]{\par\textcolor{err}{\textbf{Ошибка:}} #1\par}
\newcommand{\idea}[1]{\par\textcolor{darkaccent}{\textbf{Ключевая идея:}} #1\par}
\newcommand{\result}[1]{\par\textcolor{darkaccent}{\textbf{Итог:}} #1\par}
\newcommand{\practice}[1]{\par\medskip{\color{accent}\hrule}\smallskip\textcolor{darkaccent}{\textbf{Задача на закрепление}}\par #1}
\newcommand{\missing}{\sub{Ваше решение и ответ}На представленных фотографиях решения и ответа нет. Последний правильный шаг установить невозможно.}
\begin{document}
\begin{titlepage}
\centering
\vspace*{27mm}
{\Large\color{darkaccent} МАТЕМАТИКА\par}
\vspace{17mm}
{\Huge Ревью домашней работы\par}
\vspace{10mm}
{\LARGE\color{accent} ОГЭ. Вариант 21\par}
\vspace{17mm}
{\large Разбор первой части\par}
\vfill
{\large Ярослав Верещагин\par}
\vspace{8mm}
26 ноября 2026 г.\par
\vspace{18mm}
{\small\signature\par}
\vspace{17mm}
\begin{tikzpicture}
\draw[accent,line width=.8pt] (0,0) .. controls (2,0.65) and (4,-0.65) .. (6,0) .. controls (8,0.65) and (10,-0.65) .. (12,0);
\end{tikzpicture}
\end{titlepage}
\setcounter{page}{1}
{\Large\color{darkaccent} Сводная таблица}\par\medskip
В таблицу включены только пропущенные задания первой части и решение, в котором требуется исправить ход рассуждения.
\par\medskip
\renewcommand{\arraystretch}{1.45}
\begin{tabularx}{\linewidth}{@{}p{.07\linewidth}p{.20\linewidth}X p{.15\linewidth}p{.15\linewidth}@{}}
\toprule
№ & Тема & Суть замечания & Ваш ответ & Верный ответ\\
\midrule
\hyperref[task:5]{5} & Масштаб & Нет решения & --- & 14\\
\hyperref[task:13]{13} & Неравенство & Не обоснован переход к интервалу & $(-\infty;-1{,}5)$ & 2\\
\hyperref[task:14]{14} & Степени & Нет решения & --- & 2,5\\
\hyperref[task:15]{15} & Треугольник & Нет решения & --- & 21\\
\hyperref[task:17]{17} & Трапеция & Нет решения & --- & 133\\
\hyperref[task:18]{18} & Площадь & Нет решения & --- & 20\\
\bottomrule
\end{tabularx}
\clearpage
\vspace*{\fill}
\begin{center}
{\Huge\color{darkaccent} Разбор ошибок}\par\bigskip
{\color{accent}\rule{.55\linewidth}{.5pt}}
\end{center}
\vspace*{\fill}
\task{5}{Размер шрифта}
\sub{Условие}
Текст напечатан на листе A4 шрифтом высотой 10 пунктов. Найдите высоту шрифта для листа A3 при сохранении пропорций текста. Ответ округлите до целого числа пунктов. Размеры A4: $210\times297$ мм, A3: $297\times420$ мм. Один пункт равен $0{,}3528$ мм.
\missing
\error{Пропуск задания.} Необходимо установить линейный коэффициент увеличения, а не отношение площадей листов.
\idea{При подобии все линейные размеры изменяются в одно и то же число раз. Высота букв является линейной величиной.}
\sub{Полное решение}
Лист A3 получается из A4 увеличением ширины примерно в $297/210$ раза. Обозначим этот коэффициент через $k$:
\[
k=\frac{297}{210}\approx1{,}4143.
\]
Длина увеличивается практически в то же число раз: $420/297\approx1{,}4141$. Незначительное различие объясняется округлением размеров листов до миллиметров.
Высоту шрифта необходимо умножить на тот же коэффициент:
\[
h=10\cdot\frac{297}{210}\approx14{,}1429\text{ пт}.
\]
Переводить пункты в миллиметры необязательно: при вычислении отношения одинаковые единицы сокращаются. При округлении до целых получаем 14.
\sub{Проверка}
Площадь листа увеличивается вдвое, поэтому линейный размер увеличивается примерно в $\sqrt{2}$ раза, а не в два раза. $10\sqrt{2}\approx14{,}14$, что подтверждает результат.
\result{Ответ: $14$.}
\practice{На листе A4 текст напечатан шрифтом высотой 12 пунктов. Найдите высоту шрифта на листе A3 при сохранении пропорций. Размеры листов соответственно $210\times297$ мм и $297\times420$ мм. Ответ округлите до целого.}
\task{13}{Линейное неравенство}
\sub{Условие}
Укажите решение неравенства $6x-3(4x+1)>6$.
\sub{Ваше решение и ответ}
Сначала правильно раскрыты скобки: $6x-12x-3>6$. Далее записано равенство $6x-12x-3-6=0$; найдена граница $x=-1{,}5$ и указан интервал $(-\infty;-1{,}5)$.
\error{Не обоснован переход от неравенства к уравнению и выбор интервала.} Само уравнение можно использовать для поиска границы, но оно не равносильно исходному неравенству. Верный конечный интервал сохраняем.
\idea{При делении обеих частей неравенства на отрицательное число знак неравенства меняется на противоположный.}
\sub{Исправление и полное решение}
Продолжим с последнего правильного шага. Переносим свободное слагаемое вправо и приводим подобные:
\begin{align*}
6x-12x-3&>6,\\
-6x&>9,\\
x&<\frac{9}{-6}=-1{,}5.
\end{align*}
При делении на $-6$ знак $>$ заменяется на $<$. Граница не входит в множество решений, поскольку неравенство строгое.
\sub{Проверка}
При $x=-2$ левая часть равна $-12-3(-8+1)=9>6$. При $x=0$ она равна $-3$, что меньше 6. При $x=-1{,}5$ левая часть равна 6, поэтому точка исключается. В вариантах ответа найденному интервалу соответствует номер 2.
\result{$x\in(-\infty;-1{,}5)$. Ответ: $2$.}
\practice{Решите неравенство $4x-2(5x+3)>12$. Запишите множество решений в виде числового промежутка.}
\task{14}{Период полураспада}
\sub{Условие}
Масса радиоактивного изотопа уменьшается вдвое каждые 9 минут. Начальная масса равна 320 мг. Найдите массу через 63 минуты. Ответ дайте в миллиграммах.
\missing
\error{Пропуск задания.} Здесь нужно считать число полных периодов, а затем делить массу пополам за каждый период.
\idea{При каждом очередном периоде половина берётся от оставшейся массы, а не от начальной. После $n$ периодов масса равна $m_0/2^n$.}
\sub{Полное решение}
Сначала найдём количество периодов полураспада:
\[
n=\frac{63}{9}=7.
\]
Следовательно, массу необходимо семь раз разделить на два. Поскольку $2^7=128$, получаем
\[
m=\frac{320}{2^7}=\frac{320}{128}=2{,}5\text{ мг}.
\]
\sub{Проверка}
Выполним последовательное деление пополам. Каждая стрелка соответствует девяти минутам:
\[
320\to160\to80\to40\to20\to10\to5\to2{,}5.
\]
Семь стрелок соответствуют 63 минутам. Кроме того, $2{,}5\cdot128=320$, что подтверждает расчёт.
\result{Ответ: $2{,}5$.}
\practice{Масса изотопа уменьшается вдвое каждые 8 минут. Начальная масса составляет 480 мг. Найдите массу через 48 минут. Ответ дайте в миллиграммах.}
\task{15}{Биссектриса равностороннего треугольника}
\sub{Условие}
Сторона равностороннего треугольника равна $14\sqrt{3}$. Найдите биссектрису этого треугольника.
\missing
\error{Пропуск задания.} Необходимо использовать свойство биссектрисы равностороннего треугольника и теорему Пифагора.
\idea{В равностороннем треугольнике биссектриса одновременно является медианой и высотой.}
\sub{Полное решение}
Пусть $ABC$ --- равносторонний треугольник, $AH$ --- биссектриса угла $A$, точка $H$ лежит на $BC$. Так как $AB=AC$ и углы $BAH$ и $HAC$ равны, треугольники $ABH$ и $ACH$ равны по двум сторонам и углу между ними. Значит, $BH=HC$, а равные смежные углы при $H$ составляют по $90^\circ$.
\[
BH=\frac{BC}{2}=7\sqrt{3}.
\]
В прямоугольном треугольнике $ABH$ гипотенуза равна $AB=14\sqrt{3}$. По теореме Пифагора:
\begin{align*}
AH^2&=AB^2-BH^2\\
&=(14\sqrt{3})^2-(7\sqrt{3})^2\\
&=588-147=441.
\end{align*}
Длина положительна, поэтому $AH=\sqrt{441}=21$.
\sub{Проверка}
По формуле высоты равностороннего треугольника $h=a\sqrt{3}/2$ получаем $14\sqrt{3}\cdot\sqrt{3}/2=21$. Результат совпадает.
\result{Ответ: $21$.}
\practice{Сторона равностороннего треугольника равна $10\sqrt{3}$. Найдите длину его биссектрисы.}
\task{17}{Углы равнобедренной трапеции}
\sub{Условие}
Сумма двух углов равнобедренной трапеции равна $94^\circ$. Найдите больший угол этой трапеции. Ответ дайте в градусах.
\missing
\error{Пропуск задания.} Требуется определить, какие два угла могли дать указанную сумму.
\idea{Углы при каждом основании равнобедренной трапеции равны, а углы при одной боковой стороне в сумме составляют $180^\circ$.}
\sub{Полное решение}
Обозначим меньший угол через $\alpha$, больший через $\beta$. Тогда углы трапеции имеют величины $\alpha,\alpha,\beta,\beta$, причём
\[
\alpha+\beta=180^\circ.
\]
Два угла при разных основаниях не могут давать $94^\circ$, поскольку их сумма равна $180^\circ$. Два больших угла тоже не подходят: каждый из них больше $90^\circ$. Следовательно, $94^\circ$ --- сумма двух меньших углов.
\[
2\alpha=94^\circ,\qquad \alpha=47^\circ.
\]
Больший угол дополняет меньший до развёрнутого:
\[
\beta=180^\circ-47^\circ=133^\circ.
\]
\sub{Проверка}
Сумма двух меньших углов равна $47^\circ+47^\circ=94^\circ$. Сумма всех углов равна $2\cdot47^\circ+2\cdot133^\circ=360^\circ$. Условия выполнены.
\result{Ответ: $133$.}
\practice{Сумма двух углов равнобедренной трапеции равна $112^\circ$. Найдите больший угол трапеции. Ответ дайте в градусах.}
\task{18}{Площадь параллелограмма}
\sub{Условие}
На клетчатой бумаге с размером клетки $1\times1$ изображён параллелограмм. Найдите его площадь. На рисунке варианта вертикальная сторона занимает 5 клеток, расстояние между вертикальными сторонами --- 4 клетки.
\missing
\error{Пропуск задания.} Для площади нужно выбрать сторону и перпендикулярную ей высоту, а не умножать длины двух соседних сторон.
\idea{Площадь параллелограмма равна произведению основания на соответствующую высоту: $S=ah$. Основанием может быть любая сторона.}
\sub{Полное решение}
Удобно выбрать основанием вертикальную сторону. Она имеет длину 5: считаем промежутки между горизонтальными линиями сетки, а не сами линии. Высота перпендикулярна вертикальной стороне, поэтому направлена горизонтально. Между вертикальными сторонами ровно 4 клетки.
\begin{center}
\begin{tikzpicture}[scale=.55]
\draw[step=1,gray!30,very thin] (-1,-1) grid (5,9);
\draw[accent,thick] (0,3)--(0,8)--(4,5)--(4,0)--cycle;
\draw[darkaccent,dashed,thick] (0,4)--(4,4) node[midway,above,fill=white] {$h=4$};
\draw[darkaccent] (0,4.35)--(.35,4.35)--(.35,4);
\node[left,fill=white] at (0,5.5) {$a=5$};
\end{tikzpicture}
\end{center}
\[
S=5\cdot4=20.
\]
\sub{Проверка}
Параллелограмм вписан в прямоугольник шириной 4 и высотой 8. Вне фигуры остаются два прямоугольных треугольника с катетами 4 и 3. Поэтому $4\cdot8-2\cdot(4\cdot3/2)=20$.
\result{Ответ: $20$.}
\practice{На клетчатой бумаге с клеткой $1\times1$ изображён параллелограмм с вершинами $(1;1)$, $(1;7)$, $(5;9)$ и $(5;3)$. Найдите его площадь.}
\vfill
\begin{center}\small\signature\end{center}
\end{document}
